% ============================================================================
% CAPITULO 1: INTRODUCCION
% ============================================================================

\section{Introducci\'on}
\label{sec:introduccion}

% TODO: Redactar contenido usando skill tfg-cap1-introduccion

% ----------------------------------------------------------------------------
% 1.1 Contexto y Motivacion
% ----------------------------------------------------------------------------
\subsection{Contexto y Motivaci\'on}
\label{subsec:contexto-motivacion}

% Contexto social: ONGs + tecnologia en Espana
% Asociacion DNI (Damos Nuestra Ilusion) en Valencia
% Necesidad de comunicacion eficiente con voluntarios
% Brecha digital en el tercer sector

% ----------------------------------------------------------------------------
% 1.2 Definicion del Problema
% ----------------------------------------------------------------------------
\subsection{Definici\'on del Problema}
\label{subsec:definicion-problema}

% Problema concreto: DNI Valencia necesita atender consultas de voluntarios
% Limitaciones del sistema actual (manual, no escalable)
% Requisitos: multicanal (web + Telegram), persistencia, privacidad

% ----------------------------------------------------------------------------
% 1.3 Objetivos
% ----------------------------------------------------------------------------
\subsection{Objetivos}
\label{subsec:objetivos}

% Objetivo general + 4-5 objetivos especificos SMART
% OG: Desarrollar un chatbot conversacional con RAG para DNI Valencia
% OE1: Implementar pipeline RAG con busqueda hibrida
% OE2: Desarrollar sistema de contexto conversacional multi-turn
% OE3: Integrar bot de Telegram con persistencia cross-sesion
% OE4: Evaluar rendimiento con framework RAGAs (>0.85 score)
% OE5: Garantizar privacidad mediante despliegue local (Ollama)

% ----------------------------------------------------------------------------
% 1.4 Alcance y Limitaciones
% ----------------------------------------------------------------------------
\subsection{Alcance y Limitaciones}
\label{subsec:alcance-limitaciones}

% Alcance: corpus DNI (16 docs, 263 chunks), web + Telegram
% Limitaciones: idioma espanol, dominio DNI, servidor UPV Ollama
% Fuera de alcance: WhatsApp, multimodalidad, fine-tuning

% ----------------------------------------------------------------------------
% 1.5 Estructura del Documento
% ----------------------------------------------------------------------------
\subsection{Estructura del Documento}
\label{subsec:estructura-documento}

% Resumen de cada capitulo (2-3 lineas por capitulo)
% Cap 2: Marco Teorico - fundamentos RAG, LLMs, evaluacion
% Cap 3: Metodologia - arquitectura, implementacion, fases
% Cap 4: Resultados - benchmarks, metricas, comparativas
% Cap 5: Conclusiones - logros, limitaciones, trabajo futuro
